\begin{proof}[Proof of \cref{obj:thm:balanced-reduction}]
Throughout, \(p=1/2\). We use the staircase notation of \cref{obj:synth_4,obj:synth_8,obj:def:staircase-saddle}.

\begin{enumerate}
\item Define the binary privacy contraction and the corresponding information value by
\[
\kappa_\varepsilon
:=\frac{e^\varepsilon-1}{e^\varepsilon+1},
\qquad
J_{\mathrm{bin}}
:=\frac{\kappa_\varepsilon^2}
{1-\kappa_\varepsilon^2\tau^2}.
\]
Since \(\varepsilon>0\), we have \(e^\varepsilon>1\), and hence
\[
0<\kappa_\varepsilon<1.
\]
The interior-mean inequalities give
\[
-1<\mu_1-\mu_0<1,
\qquad |\tau|<1.
\]
Consequently,
\[
0\le \kappa_\varepsilon^2\tau^2
\le \kappa_\varepsilon^2<1,
\qquad
1-\kappa_\varepsilon^2\tau^2>0,
\qquad
J_{\mathrm{bin}}>0.
\]
% lean: CausalSmith.Stat.LdpAteEfficiencySurface.sharp_information_pos

\item We first establish positivity of the optimized information. For every staircase pattern,
\[
1\le b_S(j)\le e^\varepsilon,
\qquad
\sum_{j=0}^3\pi_{\theta,j}=1,
\]
so
\[
1\le h_S(\theta)
=\sum_{j=0}^3\pi_{\theta,j}b_S(j)
\le e^\varepsilon.
\]
Define the singleton design by
\[
\alpha^{\mathrm{sing}}_S
:=
\begin{cases}
(e^\varepsilon+3)^{-1},& |S|=1,\\
0,& |S|\ne1,
\end{cases}
\qquad S\in\mathcal S.
\]
For each record \(j\), exactly one singleton pattern contributes
\(e^\varepsilon\) and the other three contribute \(1\). Thus
\[
\sum_{S\in\mathcal S}
\alpha^{\mathrm{sing}}_S b_S(j)
=\frac{e^\varepsilon+3}{e^\varepsilon+3}=1,
\qquad
\alpha^{\mathrm{sing}}\in\mathcal A_\varepsilon.
\]
% lean: CausalSmith.Stat.LdpAteEfficiencySurface.staircaseFeasible_singleton_design

For completeness, its objective attains a minimum. Define
\[
\begin{aligned}
a_{\mathrm{sing}}
&:=2\sum_{S\in\mathcal S}
\alpha^{\mathrm{sing}}_S
\frac{(g_S^\top v(0))(g_{S,0}+g_{S,1})}
{h_S(\theta)},\\
b_{\mathrm{sing}}
&:=\sum_{S\in\mathcal S}
\alpha^{\mathrm{sing}}_S
\frac{(g_{S,0}+g_{S,1})^2}{h_S(\theta)}
\ge0.
\end{aligned}
\]
Expanding \(v(t)=(t,t+1)^\top\) gives
\[
F_\theta(\alpha^{\mathrm{sing}},t)
=
F_\theta(\alpha^{\mathrm{sing}},0)
+a_{\mathrm{sing}}t+b_{\mathrm{sing}}t^2,
\qquad t\in\mathbb R.
\]
If \(b_{\mathrm{sing}}=0\), nonnegativity of the objective forces
\(a_{\mathrm{sing}}=0\): otherwise, at
\[
t_{\mathrm{bad}}
:=-\frac{F_\theta(\alpha^{\mathrm{sing}},0)+1}
{a_{\mathrm{sing}}},
\]
the objective would equal \(-1\). In this case \(t=0\) is a minimizer.
If \(b_{\mathrm{sing}}>0\), define
\[
t_{\mathrm{sing}}
:=-\frac{a_{\mathrm{sing}}}{2b_{\mathrm{sing}}}.
\]
Then
\[
F_\theta(\alpha^{\mathrm{sing}},t)
-F_\theta(\alpha^{\mathrm{sing}},t_{\mathrm{sing}})
=b_{\mathrm{sing}}(t-t_{\mathrm{sing}})^2\ge0.
\]
In the constant case, set \(t_{\mathrm{sing}}:=0\).
% lean: CausalSmith.Stat.LdpAteEfficiencySurface.informationObjective_exists_minimizer

The contributions from \(S=\{0\}\) and \(S=\{3\}\) satisfy
\[
\begin{aligned}
0\le
\frac{(e^\varepsilon-1)^2t_{\mathrm{sing}}^2}
{4(e^\varepsilon+3)h_{\{0\}}(\theta)}
&\le F_\theta(\alpha^{\mathrm{sing}},t_{\mathrm{sing}}),\\
0\le
\frac{(e^\varepsilon-1)^2(t_{\mathrm{sing}}+1)^2}
{4(e^\varepsilon+3)h_{\{3\}}(\theta)}
&\le F_\theta(\alpha^{\mathrm{sing}},t_{\mathrm{sing}}).
\end{aligned}
\]
Both coefficients are strictly positive. A zero minimum would therefore
force \(t_{\mathrm{sing}}=0\) and \(t_{\mathrm{sing}}+1=0\), a contradiction.
By the definition of \(J^*\),
\[
0<
\min_{t\in\mathbb R}F_\theta(\alpha^{\mathrm{sing}},t)
\le J^*(\theta,1/2,\varepsilon).
\]
% lean: CausalSmith.Stat.LdpAteEfficiencySurface.Jstar_pos_interior

\item We next bound the balanced objective. Fix
\(\alpha\in\mathcal A_\varepsilon\), and write its staircase channel as
\[
Q_\alpha(x_j,\{S\})
:=\alpha_S b_S(j),
\qquad j\in\{0,1,2,3\},\quad S\in\mathcal S.
\]
The feasibility equations make its rows probability measures. Since
\(b_S(j)\le e^\varepsilon b_S(j')\), summing the atomic inequalities gives
\[
Q_\alpha(x_j,A)
\le e^\varepsilon Q_\alpha(x_{j'},A),
\qquad A\subseteq\mathcal S.
\]
Define the grouped binary channel by
\[
\begin{aligned}
\overline Q_\alpha(+1,A)
&:=\frac{Q_\alpha(x_0,A)+Q_\alpha(x_3,A)}2,\\
\overline Q_\alpha(-1,A)
&:=\frac{Q_\alpha(x_1,A)+Q_\alpha(x_2,A)}2.
\end{aligned}
\]
Its rows are probability measures. Pairing the original row inequalities
gives
\[
\begin{aligned}
\overline Q_\alpha(+1,A)
&\le \frac{e^\varepsilon}{2}
\{Q_\alpha(x_1,A)+Q_\alpha(x_2,A)\}
=e^\varepsilon\overline Q_\alpha(-1,A),\\
\overline Q_\alpha(-1,A)
&\le e^\varepsilon\overline Q_\alpha(+1,A).
\end{aligned}
\]
The inequalities with identical binary inputs follow from
\(1\le e^\varepsilon\).
% lean: CausalSmith.Stat.LdpAteEfficiencySurface.BalancedBinary.groupedChannel_private

Introduce the reference measure and its nonnegative row densities:
\[
\nu_\alpha
:=\overline Q_\alpha(+1,\cdot)
+\overline Q_\alpha(-1,\cdot),
\qquad
f_{\alpha,\pm}
:=\frac{d\overline Q_\alpha(\pm1,\cdot)}{d\nu_\alpha}.
\]
Define, on the entire output space,
\[
a_\alpha(z)
:=\frac{f_{\alpha,+}(z)+f_{\alpha,-}(z)}2,
\qquad
\delta_\alpha(z)
:=
\begin{cases}
\displaystyle
\frac{f_{\alpha,+}(z)-f_{\alpha,-}(z)}
{f_{\alpha,+}(z)+f_{\alpha,-}(z)},
& f_{\alpha,+}(z)+f_{\alpha,-}(z)>0,\\[6pt]
0,& f_{\alpha,+}(z)+f_{\alpha,-}(z)=0.
\end{cases}
\]
Row normalization yields
\[
\int a_\alpha\,d\nu_\alpha=1,
\qquad
a_\alpha\delta_\alpha
=\frac{f_{\alpha,+}-f_{\alpha,-}}2,
\qquad
\int a_\alpha\delta_\alpha\,d\nu_\alpha=0.
\]
Privacy implies, almost everywhere,
\[
f_{\alpha,+}\le e^\varepsilon f_{\alpha,-},
\qquad
f_{\alpha,-}\le e^\varepsilon f_{\alpha,+}.
\]
At a point where their sum is positive, these inequalities imply
\[
(e^\varepsilon+1)|f_{\alpha,+}-f_{\alpha,-}|
\le
(e^\varepsilon-1)(f_{\alpha,+}+f_{\alpha,-}),
\]
and therefore \(|\delta_\alpha|\le\kappa_\varepsilon\).
Where the sum vanishes, both densities vanish and the same bound follows
from \(\delta_\alpha=0\).
% lean: Causalean.Stat.Privacy.Binary.ae_abs_contrast_le

The grouped mixture density and its derivative with respect to the
binary sign parameter are
\[
m_{\alpha,\tau}
:=\frac{1+\tau}{2}f_{\alpha,+}
+\frac{1-\tau}{2}f_{\alpha,-}
=a_\alpha(1+\tau\delta_\alpha),
\qquad
\partial_\tau m_{\alpha,\tau}=a_\alpha\delta_\alpha.
\]
Thus its Fisher information is
\[
\mathcal I_\alpha
:=\int
\frac{a_\alpha\delta_\alpha^2}
{1+\tau\delta_\alpha}\,d\nu_\alpha.
\]
Zero-density outputs contribute zero. Since
\(|\tau|<1\) and \(|\delta_\alpha|\le\kappa_\varepsilon<1\), the remaining
denominators are positive. The pointwise inequality needed below follows
from the identity
\[
J_{\mathrm{bin}}(1-\tau\delta_\alpha)
-\frac{\delta_\alpha^2}{1+\tau\delta_\alpha}
=
\frac{\kappa_\varepsilon^2-\delta_\alpha^2}
{(1-\kappa_\varepsilon^2\tau^2)(1+\tau\delta_\alpha)}
\ge0.
\]
This identity also applies when \(\tau=0\).
Multiplying by \(a_\alpha\ge0\) and integrating gives
\[
\begin{aligned}
\mathcal I_\alpha
&\le
J_{\mathrm{bin}}
\int(a_\alpha-\tau a_\alpha\delta_\alpha)\,d\nu_\alpha\\
&=J_{\mathrm{bin}}.
\end{aligned}
\]
The bound \(|a_\alpha\delta_\alpha|\le\kappa_\varepsilon a_\alpha\)
ensures the integrability of the displayed upper bound.
% lean: Causalean.Stat.Privacy.Binary.information_le_sharp

We now identify this information with the staircase objective. Balance
gives
\[
\tau=1-2\mu_0,
\qquad
\pi_\theta
=\frac12(1-\mu_0,\mu_0,\mu_0,1-\mu_0)^\top.
\]
Hence, for every \(S\in\mathcal S\),
\[
\begin{aligned}
h_S(\theta)
&=\frac{(1-\mu_0)\{b_S(0)+b_S(3)\}
+\mu_0\{b_S(1)+b_S(2)\}}2,\\
g_S^\top v(-1/2)
&=\frac{b_S(0)+b_S(3)-b_S(1)-b_S(2)}4.
\end{aligned}
\]
Define the grouped atomic probabilities by
\[
q_{\alpha,+}(S)
:=\frac{\alpha_S\{b_S(0)+b_S(3)\}}2,
\qquad
q_{\alpha,-}(S)
:=\frac{\alpha_S\{b_S(1)+b_S(2)\}}2.
\]
The discrete Fisher-information formula therefore gives
\[
\begin{aligned}
\mathcal I_\alpha
&=
\sum_{S\in\mathcal S}
\frac{\bigl\{(q_{\alpha,+}(S)-q_{\alpha,-}(S))/2\bigr\}^2}
{\{(1+\tau)q_{\alpha,+}(S)
+(1-\tau)q_{\alpha,-}(S)\}/2}\\
&=
\sum_{S\in\mathcal S}
\alpha_S\frac{(g_S^\top v(-1/2))^2}{h_S(\theta)}
=F_\theta(\alpha,-1/2).
\end{aligned}
\]
Here a summand with \(\alpha_S=0\) is defined as zero; both grouped
probabilities then vanish. When \(\alpha_S>0\), its mixture probability
is \(\alpha_Sh_S(\theta)>0\), and the displayed equality follows by
cancellation. We have proved
\[
F_\theta(\alpha,-1/2)\le J_{\mathrm{bin}}
\qquad\text{for every }\alpha\in\mathcal A_\varepsilon.
\]
% lean: CausalSmith.Stat.LdpAteEfficiencySurface.BalancedBinary.informationObjective_balanced_le_sharp

\item The feasible set is nonempty by the singleton construction above.
It is closed, and its row constraints imply
\[
0\le\alpha_S\le\sum_{S'\in\mathcal S}\alpha_{S'}\le1.
\]
Thus it is compact. Since the objective at a fixed profiling coordinate
is continuous and linear in the weights, choose
\[
\alpha^{\mathrm{env}}
\in\operatorname*{arg\,max}_{\alpha\in\mathcal A_\varepsilon}
F_\theta(\alpha,-1/2).
\]
For each feasible \(\alpha\), its profiled minimum is at most its value
at \(-1/2\). Consequently,
\[
\begin{aligned}
J^*(\theta,1/2,\varepsilon)
&\le
\max_{\alpha\in\mathcal A_\varepsilon}
F_\theta(\alpha,-1/2)\\
&=F_\theta(\alpha^{\mathrm{env}},-1/2)
\le J_{\mathrm{bin}}.
\end{aligned}
\]
Both information values are positive, so taking reciprocals yields
\[
J_{\mathrm{bin}}^{-1}
\le V^*(\theta,1/2,\varepsilon).
\]
% lean: CausalSmith.Stat.LdpAteEfficiencySurface.BalancedBinary.Jstar_balanced_le_sharp

\item We establish the channel properties needed for the matching
bound. For the ordered record alphabet in \cref{obj:synth_8}, define
\[
(\sigma_0,\sigma_1,\sigma_2,\sigma_3):=(1,-1,-1,1).
\]
These are the values of \((2W_i-1)(2Y_i-1)\).
The randomized-response channel has atomic probabilities
\[
Q(x_j,\{z\})
=
\begin{cases}
\displaystyle\frac{e^\varepsilon}{e^\varepsilon+1},
&z=\sigma_j,\\[6pt]
\displaystyle\frac1{e^\varepsilon+1},
&z=-\sigma_j,
\end{cases}
\qquad z\in\{-1,1\}.
\]
They are positive and satisfy
\[
\sum_{z\in\{-1,1\}}Q(x_j,\{z\})
=\frac{e^\varepsilon+1}{e^\varepsilon+1}=1.
\]
Thus \(Q\) is a Markov channel, and using this same kernel at each
release makes it stationary. For every pair of rows and every output,
\[
\frac{Q(x_j,\{z\})}{Q(x_{j'},\{z\})}
\in\{e^{-\varepsilon},1,e^\varepsilon\},
\]
so
\[
Q(x_j,\{z\})\le e^\varepsilon Q(x_{j'},\{z\}).
\]
Summing over the outputs in any measurable event gives
\[
Q(x_j,A)\le e^\varepsilon Q(x_{j'},A).
\]
% lean: CausalSmith.Stat.LdpAteEfficiencySurface.balancedRR_stationary

\item Group this channel in the same way:
\[
\overline Q_{\mathrm{RR}}(+1,\cdot)
:=\frac{Q(x_0,\cdot)+Q(x_3,\cdot)}2=Q(x_0,\cdot),
\qquad
\overline Q_{\mathrm{RR}}(-1,\cdot)
:=\frac{Q(x_1,\cdot)+Q(x_2,\cdot)}2=Q(x_1,\cdot).
\]
The equalities hold because the paired records have the same input
sign. Its binary-input rows are therefore ordinary randomized response.
Mixing them with probabilities \((1+\tau)/2\) and \((1-\tau)/2\) gives
\[
P_\theta(Z=z)=\frac{1+z\kappa_\varepsilon\tau}{2},
\qquad
\partial_\tau P_\theta(Z=z)=\frac{z\kappa_\varepsilon}{2},
\qquad z\in\{-1,1\}.
\]
The two mixture probabilities are positive. Its scalar Fisher
information is consequently
\[
\begin{aligned}
\mathcal I_{\mathrm{RR}}
&:=
\sum_{z\in\{-1,1\}}
\frac{\{\partial_\tau P_\theta(Z=z)\}^2}{P_\theta(Z=z)}\\
&=
\frac{\kappa_\varepsilon^2}{2(1+\kappa_\varepsilon\tau)}
+\frac{\kappa_\varepsilon^2}{2(1-\kappa_\varepsilon\tau)}
=J_{\mathrm{bin}}.
\end{aligned}
\]
% lean: CausalSmith.Stat.LdpAteEfficiencySurface.grouped_balancedRR_information_eq

To relate this scalar information to the full information matrix, define
\[
D_{\mathrm{RR}}
:=\sum_{j=0}^3Q(x_j,\cdot),
\qquad
f_{\mathrm{RR},j}
:=\frac{dQ(x_j,\cdot)}{dD_{\mathrm{RR}}},
\qquad
m_{\mathrm{RR},\theta}
:=\sum_{j=0}^3\pi_{\theta,j}f_{\mathrm{RR},j}.
\]
The two-coordinate density derivative at fixed assignment probability is
\[
\dot m_{\mathrm{RR},\theta}
:=
\begin{pmatrix}
(f_{\mathrm{RR},1}-f_{\mathrm{RR},0})/2\\
(f_{\mathrm{RR},3}-f_{\mathrm{RR},2})/2
\end{pmatrix}.
\]
Under balance,
\[
m_{\mathrm{RR},\theta}
=
\frac{(1-\mu_0)(f_{\mathrm{RR},0}+f_{\mathrm{RR},3})
+\mu_0(f_{\mathrm{RR},1}+f_{\mathrm{RR},2})}{2},
\]
and
\[
v(-1/2)^\top\dot m_{\mathrm{RR},\theta}
=
\frac{f_{\mathrm{RR},0}+f_{\mathrm{RR},3}
-f_{\mathrm{RR},1}-f_{\mathrm{RR},2}}4.
\]
The grouped reference measure and its row densities are
\[
\nu_{\mathrm{RR}}:=\frac12D_{\mathrm{RR}},
\qquad
\overline f_{\mathrm{RR},+}
:=f_{\mathrm{RR},0}+f_{\mathrm{RR},3},
\qquad
\overline f_{\mathrm{RR},-}
:=f_{\mathrm{RR},1}+f_{\mathrm{RR},2}.
\]
Its mixture density satisfies
\[
\overline m_{\mathrm{RR},\tau}
:=\frac{1+\tau}{2}\overline f_{\mathrm{RR},+}
+\frac{1-\tau}{2}\overline f_{\mathrm{RR},-}
=2m_{\mathrm{RR},\theta}.
\]
Therefore, writing
\[
I_{\mathrm{RR}}:=I_\theta(Q),
\]
we obtain
\[
\begin{aligned}
\mathcal I_{\mathrm{RR}}
&=
\int
\frac{\{(\overline f_{\mathrm{RR},+}
-\overline f_{\mathrm{RR},-})/2\}^2}
{2m_{\mathrm{RR},\theta}}\,d\nu_{\mathrm{RR}}\\
&=
\int
\frac{\{v(-1/2)^\top\dot m_{\mathrm{RR},\theta}\}^2}
{m_{\mathrm{RR},\theta}}\,dD_{\mathrm{RR}}\\
&=v(-1/2)^\top I_{\mathrm{RR}}v(-1/2)
=J_{\mathrm{bin}}.
\end{aligned}
\]
The integrands are defined as zero at zero-density outputs. Indeed, if
the sum of the four nonnegative row densities vanishes, all four vanish.
If that sum is positive and
\(f_{\mathrm{RR},0}+f_{\mathrm{RR},3}=0\), then
\[
m_{\mathrm{RR},\theta}
=\frac{\mu_0}{2}
(f_{\mathrm{RR},1}+f_{\mathrm{RR},2})>0.
\]
If \(f_{\mathrm{RR},0}+f_{\mathrm{RR},3}>0\), then
\[
m_{\mathrm{RR},\theta}
\ge\frac{1-\mu_0}{2}
(f_{\mathrm{RR},0}+f_{\mathrm{RR},3})>0.
\]
Thus the change of reference measure and the divisions above cover all
outputs.
% lean: CausalSmith.Stat.LdpAteEfficiencySurface.BalancedBinary.grouped_information_eq_channelQuadratic

The paired-row identities also give, almost everywhere,
\[
f_{\mathrm{RR},0}=f_{\mathrm{RR},3},
\qquad
f_{\mathrm{RR},1}=f_{\mathrm{RR},2}.
\]
It follows that, for every \(t\in\mathbb R\),
\[
v(t)^\top\dot m_{\mathrm{RR},\theta}
=
\frac{t(f_{\mathrm{RR},1}-f_{\mathrm{RR},0})
+(t+1)(f_{\mathrm{RR},3}-f_{\mathrm{RR},2})}{2}
=
v(-1/2)^\top\dot m_{\mathrm{RR},\theta}.
\]
Integrating the squared directional derivative therefore gives
\[
v(t)^\top I_{\mathrm{RR}}v(t)=J_{\mathrm{bin}},
\qquad t\in\mathbb R.
\]
% lean: CausalSmith.Stat.LdpAteEfficiencySurface.balancedRR_quadratic_eq

\item We now evaluate the contrast variance using this positive
minimizing quadratic value. The matrix \(I_{\mathrm{RR}}\) is symmetric,
and for every \(\xi_{\mathrm{dir}}\in\mathbb R^2\),
\[
\xi_{\mathrm{dir}}^\top I_{\mathrm{RR}}\xi_{\mathrm{dir}}
=
\int
\frac{(\xi_{\mathrm{dir}}^\top\dot m_{\mathrm{RR},\theta})^2}
{m_{\mathrm{RR},\theta}}\,dD_{\mathrm{RR}}
\ge0.
\]
Thus it is positive semidefinite.
% lean: CausalSmith.Stat.LdpAteEfficiencySurface.channelFisherInfo_posSemidef

Define the perturbation coefficients by
\[
a_{\mathrm{geom}}
:=\sum_{k,\ell=0}^1(I_{\mathrm{RR}})_{k\ell},
\qquad
b_{\mathrm{geom}}
:=2\{(I_{\mathrm{RR}}v(-1/2))_0
+(I_{\mathrm{RR}}v(-1/2))_1\}.
\]
Symmetry and expansion give, for every \(\eta\in\mathbb R\),
\[
\begin{aligned}
0
&\le
v(-1/2+\eta)^\top I_{\mathrm{RR}}v(-1/2+\eta)
-v(-1/2)^\top I_{\mathrm{RR}}v(-1/2)\\
&=\eta b_{\mathrm{geom}}+\eta^2a_{\mathrm{geom}}.
\end{aligned}
\]
If \(b_{\mathrm{geom}}\ne0\), choose
\[
\eta_{\mathrm{test}}
:=-\frac{b_{\mathrm{geom}}}{2(|a_{\mathrm{geom}}|+1)}.
\]
Using \(a_{\mathrm{geom}}\le|a_{\mathrm{geom}}|\), the last expression
would satisfy
\[
\begin{aligned}
\eta_{\mathrm{test}}b_{\mathrm{geom}}
+\eta_{\mathrm{test}}^2a_{\mathrm{geom}}
&=
\frac{b_{\mathrm{geom}}^2
\{a_{\mathrm{geom}}-2(|a_{\mathrm{geom}}|+1)\}}
{4(|a_{\mathrm{geom}}|+1)^2}\\
&\le-\frac{b_{\mathrm{geom}}^2}
{4(|a_{\mathrm{geom}}|+1)}<0.
\end{aligned}
\]
Hence \(b_{\mathrm{geom}}=0\), so the coordinates of
\(I_{\mathrm{RR}}v(-1/2)\) sum to zero. Since
\[
J_{\mathrm{bin}}
=v(-1/2)^\top I_{\mathrm{RR}}v(-1/2)
=(I_{\mathrm{RR}}v(-1/2))_1,
\]
we obtain
\[
I_{\mathrm{RR}}v(-1/2)
=J_{\mathrm{bin}}(-1,1)^\top
=J_{\mathrm{bin}}c.
\]
Define
\[
w_{\mathrm{sol}}
:=\frac{v(-1/2)}{J_{\mathrm{bin}}}\in\mathbb R^2.
\]
Then
\[
I_{\mathrm{RR}}w_{\mathrm{sol}}=c,
\qquad
c^\top w_{\mathrm{sol}}=J_{\mathrm{bin}}^{-1},
\qquad
c\in\operatorname{range}(I_{\mathrm{RR}}).
\]
For any other solution
\(\widetilde w_{\mathrm{sol}}\in\mathbb R^2\) of
\(I_{\mathrm{RR}}\widetilde w_{\mathrm{sol}}=c\), symmetry gives
\[
c^\top\widetilde w_{\mathrm{sol}}
=(I_{\mathrm{RR}}w_{\mathrm{sol}})^\top
\widetilde w_{\mathrm{sol}}
=w_{\mathrm{sol}}^\top
I_{\mathrm{RR}}\widetilde w_{\mathrm{sol}}
=w_{\mathrm{sol}}^\top c
=J_{\mathrm{bin}}^{-1}.
\]
Thus the solution-value infimum in the contrast-variance convention is
\[
c^\top I_{\mathrm{RR}}^+c
=
\inf_{\{\widetilde w_{\mathrm{sol}}:\,
I_{\mathrm{RR}}\widetilde w_{\mathrm{sol}}=c\}}
c^\top\widetilde w_{\mathrm{sol}}
=J_{\mathrm{bin}}^{-1}.
\]
% lean: CausalSmith.Stat.LdpAteEfficiencySurface.contrastVariance_eq_reciprocal_of_minimizer
% lean: CausalSmith.Stat.LdpAteEfficiencySurface.balancedRR_variance_eq_sharp

\item Apply \cref{obj:thm:finite-oracle} to \(Q\), with the constant privacy
sequence
\[
\varepsilon_n:=\varepsilon,
\qquad n\in\mathbb N.
\]
Its hypotheses hold because \(0<1/2<1\), the means are interior, the
budget is positive, and \(Q\) satisfies the stationary privacy
inequalities established above. The oracle bound gives
\[
V^*(\theta,1/2,\varepsilon)
\le c^\top I_{\mathrm{RR}}^+c
=J_{\mathrm{bin}}^{-1}.
\]
Together with the reciprocal lower bound, this proves
\[
V^*(\theta,1/2,\varepsilon)
=J_{\mathrm{bin}}^{-1}.
\]
% lean: CausalSmith.Stat.LdpAteEfficiencySurface.Vstar_balanced_eq_sharp_reciprocal

\item Since \(\kappa_\varepsilon>0\) and
\(1-\kappa_\varepsilon^2\tau^2>0\), ordinary fraction algebra gives
\[
J_{\mathrm{bin}}^{-1}
=
\frac{1-\kappa_\varepsilon^2\tau^2}{\kappa_\varepsilon^2}
=\kappa_\varepsilon^{-2}-\tau^2.
\]
Moreover,
\[
\sinh(\varepsilon/2)>0,
\qquad
e^\varepsilon=(e^{\varepsilon/2})^2.
\]
Expanding the hyperbolic functions and multiplying numerator and
denominator by \(e^{\varepsilon/2}\) yields
\[
\frac{\cosh(\varepsilon/2)}{\sinh(\varepsilon/2)}
=
\frac{e^{\varepsilon/2}+e^{-\varepsilon/2}}
{e^{\varepsilon/2}-e^{-\varepsilon/2}}
=
\frac{e^\varepsilon+1}{e^\varepsilon-1}
=\kappa_\varepsilon^{-1}.
\]
Consequently,
\[
V^*(\theta,1/2,\varepsilon)
=
\left(\frac{\cosh(\varepsilon/2)}{\sinh(\varepsilon/2)}\right)^2
-\tau^2.
\]
% lean: CausalSmith.Stat.LdpAteEfficiencySurface.contraction_inv_eq_coth

\item The row normalization and event inequalities above establish the
claimed stationary Markov privacy property. The constructed information
solution establishes estimability, and the two variance identities give
\[
c\in\operatorname{range}(I_\theta(Q)),
\qquad
c^\top I_\theta(Q)^+c
=J_{\mathrm{bin}}^{-1}
=V^*(\theta,1/2,\varepsilon).
\]
% lean: CausalSmith.Stat.LdpAteEfficiencySurface.balanced_reduction
\end{enumerate}
\end{proof}
